\documentclass[10pt,a4paper]{amsart}
\usepackage[margin=19mm]{geometry}
\usepackage{amsmath,amssymb,mathtools}
\usepackage[scr=rsfs,cal=euler]{mathalfa}
\usepackage{xfrac}
\usepackage[unicode,hidelinks,bookmarks=false]{hyperref}
\newcommand{\ie}{i.e.\ }
\newcommand{\cf}{cf.\ }
\newcommand{\resp}{respectively\ }
\newcommand{\Subsection}{\subsection}
\providecommand{\nobreakdash}{\nobreak\hbox{-}}
\setlength{\emergencystretch}{2em}
\multlinegap0pt
\usepackage{amssymb}
\usepackage{xcolor}
\usepackage{algorithm}\usepackage{algpseudocode}
\usepackage{bm}
\usepackage[shortlabels]{enumitem}
\usepackage{dsfont}
\newtheorem{proposition}{Proposition}
\newcommand{\diam}{\mathrm{Diam}}
\newcommand{\dis}{\mathrm{d}}
\newcommand{\eps}{\varepsilon}
\newcommand{\g}{\mathrm{g}}
\newcommand{\length}{\mathrm{length}}
\newcommand{\man}{\mathrm{M}}
\newcommand{\tangent}{\mathrm{T}}
\title{Gromov-Wasserstein bound between Reeb and Mapper graphs}
\author{Ziyad Oulhaj, Mathieu Carri\`ere, Bertrand Michel}
\date{Source: arXiv:2506.02810v2 (CC BY 4.0)}
\begin{document}
\maketitle

\noindent\emph{Source and licence.} Gromov-Wasserstein bound between Reeb and Mapper graphs. Version arXiv:2506.02810v2 (CC BY 4.0), distributed by arXiv under the Creative Commons Attribution 4.0 International licence, http://creativecommons.org/licenses/by/4.0/, as recorded in the arXiv licence metadata for this exact version and shown on the abstract page https://arxiv.org/abs/2506.02810v2. Full licence notice: \texttt{licenses/arxiv-2506.02810-CC-BY-4.0.txt}.\par\smallskip

\noindent\textit{Editorial scope.} Counted unit: the article's proposition prop:dis\_equiv (source0.tex lines 1663--1675). The article's complete original proof of it is delivered with it (source0.tex lines 1678--1728, one \texttt{proof} environment of 51 lines). No mathematics is written, repaired or re-derived: the statement and the proof are the article's own lines, in the article's own notation, and no retained line is substituted or re-flowed.  Retained uncounted as the source-local context the proof needs: lines 1636--1644.  Nothing was omitted from any retained span, so the package carries no omission record.  No retained line contains a citation, so the wrapper carries no bibliography and no bibliography entry.  No retained line contains a figure environment, an image-inclusion command or drawing code, so no image is delivered and the package declares no asset dependency.  Editorial formatting only, in addition: an AMS \texttt{amsart} preamble that restates the article's own declarations and macros used by the retained lines, namely the environments \texttt{proposition} on separate counters, together with the standard packages those lines need.  The retained environments print in this document as Proposition 1 in order of appearance, whereas the article prints the same items with the numbering of its own full document.  The complete native source of the article as distributed in the arXiv e-print for this exact version is bundled unchanged as \texttt{sources/179225/source0.tex}.
\begin{proposition}\label{prop:g_equiv}
Let $U$ be an open set in $\man$, on which we are given coordinates $x(p)=(x_1,...,x_d)$. 
%
For every $p\in\man$ and for every $v\in\tangent_p\man$:

$$\lambda(p)\cdot \sqrt{\g_0(v,v)}\leq \sqrt{\g(v,v)}\leq \mu(p)\cdot \sqrt{\g_0(v,v)},$$

where $\g_0=\sum_{i}dx_i\cdot dx_i$ is the Euclidean metric in the local coordinates.
\end{proposition}

\begin{proposition}\label{prop:dis_equiv}Let $p\in\man$. For every small enough neighborhood $U$ of $p$, we have that for every $q\in U$:
%\begin{equation}
%\label{eq:gamma0}
$$\lambda_0\cdot\dis_0(p,q)\leq \dis(p,q)\leq\mu_0\cdot\dis_0(p,q),$$
%\end{equation}
where $\dis_0(p,q)$ is the Euclidean distance between the representations of $p$ and $q$ in local coordinates and where
$$\lambda_0=\inf_{r\in U}\lambda(r) \textrm{ and }
\mu_0=\sup_{r\in U}\mu(r).$$
In particular $\lambda_0$ tends to $ \lambda(p)$ and
$\mu_0$  tends to  $\mu(p)$ as $\diam(U)$ tends to $0$.
%Moreover, $$\frac{\dis(p,q)}{\dis_0(p,q)}\rightarrow 1\text{ as }q\rightarrow p.$$
% \bert{ca grince un cette convergence de $\lambda_0, \mu_0$, je te propose plutot (a verifier et peut etre à positionner plutot en appendice ?) : Moreover, it is possible to choose $\lambda_0$ and $\mu_0$ continuously in the following sense: for any $r$ small enough  there exists $\lambda_0(r)$ and $\mu_0(r)$ satisfying \eqref{eq:gamma0} for any $q \in U \cap B_0(p,r)$, such that $\lambda_0(r)$ and $\mu_0(r)$ tends to $\lambda_0$ and $\mu_0$ as $r$ tends to zero.}
\end{proposition}

\begin{proof}
Let $p\in\man$ and let $U$ be a coordinate neighborhood of $p$.
Up to shrinking $U$, we can assume that:
$$U=\{q\in\man\, ,\, \dis_0(p,q)\leq\eps\},$$
for a given, small enough $\eps>0$, and where $\dis_0$ is the distance associated to the metric $\g_0$, given by 
$$\dis_0(p,q)=\sqrt{x_1(q)^2+\dots+x_d(q)^2},$$
where $q$ is given in local coordinates by $x(q)=(x_1(q),...,x_d(q))$.

From Proposition~\ref{prop:g_equiv}, there exist continuous functions $\lambda,\mu\colon U\rightarrow(0,\infty)$ such that for every $q\in U$ and $v\in\tangent_q\man$, one has
$$\lambda(q)\cdot \sqrt{\g_0(v,v)}\leq \sqrt{\g(v,v)}\leq \mu(q)\cdot \sqrt{\g_0(v,v)}.$$
Furthermore, $\lambda(q),\mu(q)\rightarrow \lambda(p),\mu(p)$ as $q\rightarrow p$ because $\lambda$ and $\mu$ are continuous.

Now, let $\gamma\colon[0,1]\rightarrow\man$ be a piecewise-$C^{\infty}$ curve such that $\gamma(0)=p$ and $\gamma(1)=q\in U$.

\begin{itemize}
\item If $\gamma$ is a segment for $\dis_0$, then it lies in $U$ (as $\dis_0(p,q)\leq\eps$). We also have:
%
\begin{align*}
\dis_0(p,q)&=\int_0^1\sqrt{\g_0\left(\frac{d\gamma(t)}{dt},\frac{d\gamma(t)}{dt}\right)}\, dt\\
&\geq \frac{1}{\sup_{t\in[0,1]}\ \mu(\gamma(t))}\int_0^1\sqrt{\g\left(\frac{d\gamma(t)}{dt},\frac{d\gamma(t)}{dt}\right)}\, dt\\
&=\frac{\length(\gamma)}{\sup_{t\in[0,1]}\ \mu(\gamma(t))}\\
&\geq \frac{\dis(p,q)}{\sup_{t\in[0,1]}\ \mu(\gamma(t))}.
\end{align*}
%
\item If $\gamma$ lies inside $U$, then:
%
\begin{align*}
\length(\gamma)&=\int_0^1\sqrt{\g\left(\frac{d\gamma(t)}{dt},\frac{d\gamma(t)}{dt}\right)}\, dt\\
&\geq \inf_{t\in[0,1]}\ \lambda(\gamma(t)) \cdot \int_0^1\sqrt{\g_0\left(\frac{d\gamma(t)}{dt},\frac{d\gamma(t)}{dt}\right)}\, dt\\
&\geq \inf_{t\in[0,1]}\ \lambda(\gamma(t))\cdot\dis_0(p,q)
\end{align*}
%
\item If $\gamma$ leaves $U$, then there exists a smallest $t_0$ such that $\gamma(t_0)\notin U$. We have:
%
\begin{align*}
\length(\gamma)&\geq\int_0^{t_0}\sqrt{\g\left(\frac{d\gamma(t)}{dt},\frac{d\gamma(t)}{dt}\right)}\, dt\\
&\geq \inf_{t\in[0,1]}\ \lambda(\gamma(t))\cdot\int_0^{t_0}\sqrt{\g_0\left(\frac{d\gamma(t)}{dt},\frac{d\gamma(t)}{dt}\right)}\, dt\\
&\geq \inf_{t\in[0,1]}\ \lambda(\gamma(t))\cdot\eps\\
&\geq \inf_{t\in[0,1]}\ \lambda(\gamma(t))\cdot\dis_0(p,q)
\end{align*}
%\mathieu{Un dessin illustrant les 3 cas pourrait être sympa.}
Like mentioned above, $\lambda(q),\mu(q)\rightarrow \lambda(p),\mu(p)$ as $q\rightarrow p$, therefore (and by shrinking $U$ again if necessary) we can ensure that $\lambda_0=\inf_{r\in U}\lambda(r)>0$ and that $\mu_0=\sup_{r\in U}\mu(r)<\infty$.

We have proved that for every $q\in U$, one has
$$\lambda_0\cdot\dis_0(p,q)\leq \dis(p,q)\leq\mu_0\cdot\dis_0(p,q),$$
and we finally see that, as $q\rightarrow p$:
$$\lambda_0,\mu_0\rightarrow \lambda(p),\mu(p).$$
%Finally, notice that by definition of $\mu_0$ and $\lambda_0$ and since $\lambda(q),\mu(q)\rightarrow 1$ as $q\rightarrow p$, we know that $$\frac{\dis(p,q)}{\dis_0(p,q)}\rightarrow 1\text{ as }q\rightarrow p.$$
\end{itemize}

\end{proof}

\end{document}
